% Midterm report: E1 (performance) and E4 (correctness).
% Build after analysis/part1_figures.py: cd report/midterm && latexmk -pdf midterm.tex
\documentclass[11pt,a4paper]{article}
\usepackage[margin=2.1cm,vmargin=1.6cm]{geometry}
\usepackage[utf8]{inputenc}
\usepackage{microtype,booktabs,enumitem,array,tabularx,float,caption}
\usepackage[hidelinks]{hyperref}
\usepackage[compact]{titlesec}
\titleformat{\section}{\large\bfseries}{\thesection.}{0.5em}{}
\titleformat{\subsection}{\normalsize\bfseries}{\thesubsection}{0.5em}{}
\setlist{leftmargin=1.5em,itemsep=2pt,topsep=3pt}
\setlength{\parindent}{0pt}
\setlength{\parskip}{6pt}
\captionsetup{font=small,labelfont=bf,labelsep=period,skip=4pt}
\setlength{\textfloatsep}{16pt}
\setlength{\intextsep}{14pt}
\DeclareUnicodeCharacter{00D7}{\ensuremath{\times}}
\DeclareUnicodeCharacter{00B5}{\ensuremath{\mu}}
\hypersetup{pdftitle={Midterm Report: Post-Quantum KEM Comparative Study},pdfauthor={Hamdan Al-Jabri}}

\begin{document}

{\raggedright\LARGE\bfseries Midterm Report: Comparative Study of Post-Quantum Key-Exchange Mechanisms\par}
\vspace{6pt}
Hamdan Al-Jabri (Neptun QUTC51)\\
Cybersecurity Lab 1, ELTE Faculty of Informatics\\
November 2026

\section{Introduction}
Data encrypted today can be recorded and decrypted later, once a large quantum computer exists. NIST has
therefore standardised ML-KEM (FIPS 203, 2024) and selected HQC as a backup (2025), and browsers already use
the hybrid X25519MLKEM768. The available post-quantum KEMs, however, differ by orders of magnitude in speed and
size, and their implementations differ in quality.

The project compares them in four experiments: \textbf{E1} performance, \textbf{E2} real TLS 1.3 handshakes,
\textbf{E3} timing leaks, and \textbf{E4} correctness. This report covers E4 and E1. Code and data are public
at \url{https://github.com/hsh1994-ux/pq-kem-study}.

\section{Background}
\subsection{Key-encapsulation mechanisms}
A KEM lets two parties agree on a shared secret in three steps:
\begin{enumerate}
  \item \textbf{KeyGen}: the client makes a key pair and sends the public key.
  \item \textbf{Encaps}: the server uses the public key to create a random shared secret and a ciphertext, and
  sends the ciphertext back.
  \item \textbf{Decaps}: the client uses its private key to recover the same secret from the ciphertext.
\end{enumerate}

\subsection{The four methods: structured and unstructured}
Post-quantum KEMs rest on two main kinds of hard maths problems:
\begin{itemize}
  \item \textbf{Lattices:} the public key hides a secret behind a large system of linear equations with small
  random errors added (the \emph{Learning With Errors} problem). Without the secret, removing the errors is
  believed to be hard, even for a quantum computer. ML-KEM uses \emph{Module}-LWE, where the equations are
  built from polynomials; FrodoKEM uses plain LWE.
  \item \textbf{Error-correcting codes:} the public key is a disguised version of a code that the owner can
  decode efficiently. The ciphertext is a message with deliberate errors added; without the secret
  disguise, correcting those errors is believed to be hard. This idea goes back to McEliece in 1978.
\end{itemize}

Both problems need a large matrix as the public key. The difference between the methods is \textbf{how that
matrix is built}:
\begin{itemize}
  \item \textbf{Structured:} the matrix follows a repeating pattern, so it can be described by a few short
  rows instead of being stored in full. In ML-KEM the matrix entries are polynomials, and one polynomial
  stands for a whole block of the matrix; in HQC the matrix is built from circulant blocks, where each row of
  a block is a cyclic shift of the row before it. This makes keys \textbf{much smaller} and the arithmetic
  \textbf{much faster}. The price is extra algebraic structure, which adds assumptions to the security
  argument and needs its own cryptanalysis. No practical attack that uses it is known today.
  \item \textbf{Unstructured:} the public key has no repeating pattern and is stored or generated in full.
  FrodoKEM uses plain LWE on a large random matrix. Classic McEliece hides a binary Goppa code behind a
  random-looking matrix; the Goppa code itself has algebraic structure (it is the secret trapdoor), but the
  public key shows no pattern, and the design has resisted attacks for over 45 years. Keys are \textbf{much
  larger} and the work is \textbf{slower}, but these designs avoid the extra assumptions of structured ones,
  which is why they are seen as the more \emph{conservative} choice. Neither kind is automatically more secure.
\end{itemize}

Choosing one method of each kind lets the study answer two questions at once: how lattices compare with codes,
and \textbf{what the extra caution of unstructured designs costs}. Our own measurements (Section~4) show the
size of this cost: at level 3, ML-KEM sends 2,272 bytes and FrodoKEM 31,424 bytes (14× more); HQC sends 13,492
bytes and Classic McEliece 524,316 bytes (39× more, almost all of it the public key). Structure is not the
only difference, however: the methods also differ in parameters, arithmetic and how well optimised their code is.

\begin{table}[H]
\centering\small
\begin{tabular}{@{}llll@{}}
\toprule
Method & Based on & Design & Status \\
\midrule
ML-KEM & lattices & structured & NIST standard, FIPS 203 (2024) [1] \\
HQC & codes & structured & NIST backup (2025); standard due 2027 [2] \\
FrodoKEM & lattices & unstructured & ISO/IEC 18033-2 Amd 2 (2026); recommended by BSI [3,5] \\
Classic McEliece & codes & unstructured & ISO/IEC 18033-2 Amd 2 (2026); recommended by BSI [4,5] \\
\bottomrule
\end{tabular}
\caption{The four KEMs, one of each kind. Each is tested at NIST security levels 1, 3 and 5 (as strong as AES-128, -192 and
-256), so 12 versions in total.}
\end{table}

\subsection{Baselines}
\begin{itemize}
  \item \textbf{X25519 DH} [6]: plain Diffie--Hellman, as used in TLS today. It is not a KEM, but it has the same
  three steps: the client makes a key pair (KeyGen); the server makes its own key pair and computes the secret
  (Encaps); the client computes the same secret (Decaps). This is the main baseline.
  \item \textbf{X25519 KEM} (DHKEM, RFC 9180 [7]): the same calculation inside the KEM interface, plus one step
  that hashes the result with both public keys. Measured for comparison.
  \item \textbf{Hybrid X25519MLKEM768} [8]: a TLS 1.3 key-exchange group, standardised in RFC 10024 (2026) and
  used by browsers today. It is packaged as one KEM, but inside it is
  plain X25519 DH plus ML-KEM-768; the two secrets are joined and TLS hashes them later. It stays safe if
  either part is broken.
\end{itemize}

\section{Method}
\subsection{Setup}
Two computers were used: an \textbf{Apple M4} MacBook Air (ARM64, Linux in Docker) and an \textbf{Intel
i7-1185G7} laptop (x86-64, Linux in WSL2). Both run the same software: liboqs 0.16.0 [9] for the
post-quantum KEMs and OpenSSL 3.5.8 for X25519, the hybrid and a second ML-KEM implementation. Both use the
same Linux system (Ubuntu 24.04), compiler (GCC 13.3), build flags and benchmark settings. Besides the
processor, the host operating system (macOS, Windows) and the virtualisation (Docker, WSL2) also differ.

liboqs was built in two versions: \textbf{portable C}, which runs on any processor, and \textbf{optimised}
code, which uses vector instructions (AVX2 on Intel, NEON on the M4). Optimised code does not exist for every
method on every processor:

\begin{table}[H]
\centering\small
\begin{tabular}{lcc}
\toprule
Method & Intel (AVX2) & Apple M4 (NEON) \\
\midrule
ML-KEM & yes & yes \\
HQC & no & no \\
FrodoKEM & yes & no \\
Classic McEliece & yes & no \\
\bottomrule
\end{tabular}
\caption{Optimised code in liboqs 0.16.0. Where there is none, only the portable code is used.}
\end{table}

\subsection{E4: correctness}
Before measuring speed, five kinds of checks were run on both computers:
\begin{enumerate}
  \item \textbf{Known answers:} the code gives exactly the expected output for fixed inputs (all 12, both versions).
  \item \textbf{NIST test vectors:} ML-KEM matches NIST's official FIPS 203 test vectors (ML-KEM only, the
  other methods have no NIST standard yet).
  \item \textbf{Random exchanges:} with random inputs both sides always get the same secret (all 12, both versions).
  \item \textbf{Cross-version:} keys made by the portable code work with the optimised code, and the other way
  round (all 12).
  \item \textbf{Cross-library:} liboqs and OpenSSL ML-KEM work together (ML-KEM only, OpenSSL does not
  implement the other three).
\end{enumerate}

\subsection{E1: performance}
\begin{itemize}
  \item \textbf{Timing:} KeyGen, Encaps and Decaps were each run \textbf{1,000 times} and the median is
  reported. One whole key exchange = the sum of the three medians (KeyGen + Encaps + Decaps).
  \item \textbf{Classic McEliece key generation was run 100 times, not 1,000.} One run takes 40--500\,ms, so
  1,000 runs would take hours. Its time also varies a lot, because it builds a random key and starts again
  when that fails (about 7 times in 10); 100 runs give a much more stable median than 20, but its KeyGen median can still differ by up to 1.7× between runs.
  \item \textbf{Fair conditions:} the timer was calibrated at the start of every run, the processor was warmed
  up for 200\,ms before timing (the M4 needs about 90\,ms to reach full speed), and there was a 30-second
  pause between methods so the fanless MacBook could cool down. The processor's clock speed was also measured,
  so results can also be given as \emph{estimated} CPU cycles (time × measured clock speed) and compared with
  published numbers.
  \item \textbf{Also recorded:} bytes sent over the network and peak memory (stack) of each step.
  \item \textbf{Comparisons} with X25519 always use the faster of the two code versions.
\end{itemize}

\section{Results}
\subsection{E4: correctness}
All \textbf{69 checks (1,332 test cases) pass on both computers}. The optimised code gives exactly the same
results as the portable code, so comparing their speed is fair, and liboqs agrees with NIST's official
vectors and with OpenSSL's independent ML-KEM.

\subsection{E1: performance}
\begin{table}[t]
\centering\scriptsize\setlength{\tabcolsep}{3.6pt}
\begin{tabular}{lrrrrrrrr}
\toprule
 & & & \multicolumn{3}{c}{\textbf{Apple M4 (NEON code)}} & \multicolumn{3}{c}{\textbf{Intel i7 (AVX2 code)}} \\ \cmidrule(lr){4-6}\cmidrule(lr){7-9}
Method & Level & Bytes & Portable & Optimised & vs X25519 & Portable & Optimised & vs X25519  \\
\midrule
X25519 DH (today) & – & 64 & – & 62.5\,µs & 1.00× & – & 115\,µs & 1.00×  \\
X25519 KEM (DHKEM) & – & 64 & – & 71.5\,µs & 1.14× & – & 126\,µs & 1.10×  \\
Hybrid X25519MLKEM768 & – & 2,336 & – & 100\,µs & 1.60× & – & 191\,µs & 1.66×  \\
\midrule
ML-KEM-512 & 1 & 1,568 & 26.8\,µs & 17.6\,µs {\scriptsize(1.53×)} & 0.28× & 64.3\,µs & 20.8\,µs {\scriptsize(3.1×)} & 0.18×  \\
ML-KEM-768 & 3 & 2,272 & 42.7\,µs & 28.0\,µs {\scriptsize(1.53×)} & 0.45× & 98.6\,µs & 33.2\,µs {\scriptsize(3.0×)} & 0.29×  \\
ML-KEM-1024 & 5 & 3,136 & 64.5\,µs & 42.4\,µs {\scriptsize(1.52×)} & 0.68× & 143\,µs & 46.4\,µs {\scriptsize(3.1×)} & 0.40×  \\
\midrule
HQC-1 & 1 & 6,674 & 3.39\,ms & – & 54× & 7.41\,ms & – & 65×  \\
HQC-3 & 3 & 13,492 & 10.1\,ms & – & 162× & 22.3\,ms & – & 195×  \\
HQC-5 & 5 & 21,658 & 22.3\,ms & – & 357× & 52.8\,ms & – & 461×  \\
\midrule
Frodo-640 & 1 & 19,368 & 2.70\,ms & – & 43× & 6.33\,ms & 2.23\,ms {\scriptsize(2.8×)} & 19×  \\
Frodo-976 & 3 & 31,424 & 6.10\,ms & – & 98× & 13.9\,ms & 4.61\,ms {\scriptsize(3.0×)} & 40×  \\
Frodo-1344 & 5 & 43,216 & 11.0\,ms & – & 176× & 25.3\,ms & 7.58\,ms {\scriptsize(3.3×)} & 66×  \\
\midrule
McEliece-348864 & 1 & 261,216 & 65.5\,ms & – & 1,047× & 90.1\,ms & 38.7\,ms {\scriptsize(2.3×)} & 337×  \\
McEliece-460896 & 3 & 524,316 & 159\,ms & – & 2,536× & 276\,ms & 113\,ms {\scriptsize(2.5×)} & 983×  \\
McEliece-6688128 & 5 & 1,045,200 & 461\,ms & – & 7,375× & 592\,ms & 134\,ms {\scriptsize(4.4×)} & 1,166×  \\
\bottomrule
\end{tabular}

\caption{One whole key exchange: sum of the KeyGen, Encaps and Decaps medians. Portable = portable C; (\emph{n}×) after an optimised time = how many times faster it is than
portable C; -- = no optimised code. ``vs X25519'' compares with X25519 DH on the same computer, using the
faster version: below 1.00× means faster than X25519.}
\label{tab:e1}
\end{table}

Table~\ref{tab:e1} gives one whole key exchange for every method. \textbf{Note:} the time of each step
separately (KeyGen, Encaps, Decaps, both code versions) was also measured and tabulated; it is in the project repository.

\textbf{Main findings:}
\begin{itemize}
  \item \textbf{ML-KEM is faster than today's X25519} at every level: ML-KEM-768 needs 0.45× (M4) and 0.29×
  (Intel) of the X25519 time. Its cost is size: 2,272 bytes instead of 64 (35×), still small enough for one
  TLS message.
  \item \textbf{The hybrid costs 1.6--1.7× X25519}, about X25519 plus ML-KEM added together.
  The KEM form of X25519 is 10--14\% slower than plain DH.
  \item \textbf{The cautious methods are much more expensive in both time and size.} At level 1: HQC 54--65×
  X25519, FrodoKEM 19--43×, Classic McEliece 337--1,047×. McEliece's public key (0.26--1\,MB) is too large for
  a TLS key share, and it needs up to 1.5\,MiB of stack, against 10--21\,KiB for ML-KEM.
  \item \textbf{Optimised code matters as much as the method.} On Intel it makes ML-KEM and FrodoKEM about
  3× faster and McEliece Decaps 510--1,400× faster. On the M4 only ML-KEM has optimised code (1.5×). As a
  result the M4, which is about 2.2× faster on portable code, is slower than the older Intel chip for FrodoKEM
  and McEliece.
  \item \textbf{Agreement with published results:} in estimated CPU cycles, HQC and McEliece Encaps/Decaps agree
  with the designers' numbers within 7--20\%; ML-KEM and FrodoKEM are faster than older published figures
  because liboqs uses newer code.
\end{itemize}

\section{Limitations}
Both computers run Linux in a virtual machine (Docker, WSL2), and there is one computer per architecture.
Classic McEliece key generation varies strongly from run to run because of its random retries (up to 1.7×
between runs of the same code), and
McEliece-348864 key generation takes 2.6× the published cycle count, which is still being investigated.
liboqs has no optimised HQC code, so HQC's gap to ML-KEM is partly an implementation gap.

\section{Remaining work}
\textbf{E2}: real TLS 1.3 handshakes, also under network delay and packet loss, with every KEM that fits in a
TLS key share; Classic McEliece cannot be used, because its public key exceeds the 65,535-byte limit.
\textbf{E3}: tests for secret-dependent timing (dudect, Valgrind, KyberSlash-style division checks).
Then the final report and the demo at the mini-conference.



\section*{References}
\footnotesize
\begin{enumerate}[label={[\arabic*]},leftmargin=2em,itemsep=0pt]
  \item NIST, \emph{FIPS 203: Module-Lattice-Based Key-Encapsulation Mechanism Standard}, 2024.
  \item NIST, \emph{IR 8545: Status Report on the Fourth Round of the NIST PQC Standardization Process}, 2025.
  \item E.~Alkim et al., \emph{FrodoKEM specification}; ISO/IEC 18033-2:2006/Amd 2:2026.
  \item D.~J.~Bernstein et al., \emph{Classic McEliece: conservative code-based cryptography}; ISO/IEC 18033-2:2006/Amd 2:2026.
  \item BSI, \emph{TR-02102-1: Cryptographic Mechanisms: Recommendations and Key Lengths}.
  \item A.~Langley, M.~Hamburg, S.~Turner, \emph{RFC 7748: Elliptic Curves for Security}, 2016.
  \item R.~Barnes et al., \emph{RFC 9180: Hybrid Public Key Encryption}, 2022.
  \item K.~Kwiatkowski et al., \emph{RFC 10024: Post-quantum hybrid ECDHE-MLKEM key agreement for TLSv1.3}, 2026.
  \item Open Quantum Safe, \emph{liboqs 0.16.0}, \url{https://openquantumsafe.org}.
\end{enumerate}

\end{document}
